% Generado por src/comparativa.py — no editar a mano, se regenera.
\begin{table}[htbp]
    \centering
    {\footnotesize
    \begin{tabular}{@{}lrrrrrrr@{}}
        \toprule
        \textbf{Observable} & \textbf{n (cobertura)} & \textbf{Sin mitigar} & \textbf{Ridge} & \textbf{RF} & \textbf{SAGE} & \textbf{GIN} & \textbf{GATv2} \\
        \midrule
        media de $Z$ & 1986 (50\%) & 0,0046 & 0,0139 & 0,0031 & 0,0040 & 0,0048 & 0,0044 \\
        media de $X$ & 3393 (85\%) & 0,0177 & 0,0088 & 0,0034 & 0,0048 & 0,0037 & 0,0042 \\
        media de $Y$ & 1096 (27\%) & 0,0044 & 0,0109 & 0,0028 & 0,0039 & 0,0032 & 0,0034 \\
        paridad & 1825 (46\%) & 0,0111 & 0,0324 & 0,0030 & 0,0098 & 0,0062 & 0,0143 \\
        corr. de vecinos & 3369 (84\%) & 0,0072 & 0,0373 & 0,0071 & 0,0076 & 0,0058 & 0,0105 \\
        \bottomrule
    \end{tabular}}

    \caption{MAE del observable mitigado frente al de no mitigar. Cualquier valor por encima de la columna «Sin mitigar» significa que ese modelo empeora el resultado. Conjunto de test, por observable. $n$ es el número de casillas evaluables. Ridge = Ridge de 34 variables, RF = Random Forest; SAGE, GIN y GATv2 son las tres convoluciones del MPNN.  Son métricas post-ejecución: se reportan, pero la comparativa la decide el $R^2$ sobre $f$.}
    \label{tab:mitigacion_mae}

    {\footnotesize Fuente: elaboración propia.}
\end{table}
